%=====================================================================
\section*{How to read this document}
\addcontentsline{toc}{section}{How to read this document}
\begin{itemize}
  \item \textbf{Order.} Part I explains the model and every analysis method in enough detail to reproduce it. Part II is
        circles, Part III ellipses and Cassini ovals; both follow the same sequence (passive $\to$ where active defects
        sit $\to$ charge split $\to$ wall lattice $\to$ flow and tilt $\to$ survival and docking), so the same section number
        compares shapes. Part IV collects the conclusions, the status of each hypothesis and what is still missing.
  \item \textbf{Everything is in.} Nothing has been moved to a supplement yet; a grey ``SI?'' in the margin marks material
        that is a natural candidate for the supporting information. Grey boxes mark open questions.
  \item \textbf{Units.} Lengths in MPCD cells, times in MPCD time units ($\Delta t=0.1$, so 10 steps per time unit).
        $\la=22.2$ cells is the active length at the working activity $\zeta=0.07$; cavity sizes are given as $R/\la$
        (circles) or half-length$/\la$ (ovals). $\rho=2|k|/P$ is the linear density of imposed defects on a wall of
        perimeter $P$; the imposed spacing is $1/\rho$ cells.
  \item \textbf{Seeds.} Unless stated otherwise: circles 3 seeds per point (going to 10, running on the server since
        5 Oct), small ovals and the oval size series 3 seeds, ellipses 2, the largest ovals 1--3. Values quoted as a range are
        the range over seeds or over $k$, as stated.
  \item \textbf{Where the files are.} Figures and tables are copies; their sources are listed in \texttt{RESULTS.md} and
        \texttt{results/*/FILES.md}. New figures made for this document are produced by the scripts in
        \texttt{reports/master/scripts/}.
\end{itemize}

\section*{Results at a glance}
\addcontentsline{toc}{section}{Results at a glance}
\begin{table}[H]\centering\small
\caption{Status of the main claims. Confidence: \emph{high} = several seeds and independent measures agree;
\emph{good} = several seeds, one main measure; \emph{tentative} = one seed or one measure.}\label{tab:glance}
\begin{tabularx}{\textwidth}{>{\raggedright\arraybackslash}p{4.3cm}X>{\raggedright\arraybackslash}p{1.8cm}}\toprule
Claim & Evidence & Confidence\\\midrule
Passive cavities: all $2|k|$ imposed defects sit in a wall layer; bulk charge 0 & 468 passive circles; depth set by the imposed spacing, $s\approx0.56+0.094/\rho$ cells, for both signs and all $R$ (Sec.~\ref{sec:passive}) & high\\
Active circles: bulk $+1$, wall layer of $2k-2$ pinned $+\half$ & $R=1.43\,\la$, $k=2$--10: $q_b=1.00\pm0.02$; main sweep: $R\approx0.75$--$2\,\la$ at every $\rho\le0.15$; six criteria agree (Sec.~\ref{sec:ccharge}) & high\\
The $+1$ bulk is an active effect & passive bulk 0 at the same $(R,k)$ & high\\
Negative winding: bulk stays positive, reaches $+1$ only for $|k|\gtrsim15$ & 570 mirror runs; the $-\half$ wall lattice has the $+1$ period but empty sites (Sec.~\ref{sec:mirror}) & good\\
Wall layer = pinned lattice, never sliding & wall Fourier on 930 circles and 81 ovals/ellipses (Secs.~\ref{sec:clattice}, \ref{sec:slattice}) & high\\
Wall $+\half$ lean downstream by 60--65$^\circ$; circulation collapses on $R/\la$ & 930 circles (Sec.~\ref{sec:cflow}) & high\\
The $+1$ bulk is born with the vortex: circulation, tilt and bulk charge switch on together at $R/\la\approx0.4$, for every $k$ and $\rho$ & onset analysis of 151 $(\rho,k)$ points, $\rho=0.07$--0.15 (Sec.~\ref{sec:onset}) & good\\
Activity does not enter only through $\la$; $\zeta=0.07$ is the clean working point & five activities, 3 seeds (Sec.~\ref{sec:activity}) & good\\
\textbf{Elongated cavities lose the layer at fixed size and density: the split is geometric} & small ovals (3 seeds), size series, ellipses, matched in every size measure; $f^*$ collapses on the aspect ratio (Sec.~\ref{sec:scharge}) & good--high\\
Not the neck: one lobe-shaped convex oval loses it too & 3 seeds (Sec.~\ref{sec:scharge}) & good\\
Crowded walls (large $k$) restore the layer & $k$ scan (3 seeds), pinched oval $k=12$--20 (1 seed), mirror runs & good\\
Wall defects in elongated cavities escape 10--30$\times$ faster and creep upstream into empty sites & tracks, Kaplan--Meier survival, per-defect drift (Secs.~\ref{sec:slattice}, \ref{sec:ssurv}) & good\\
Two vortices per oval from aspect $\approx1.5$ & lobe rotation, 1 seed per shape (Sec.~\ref{sec:sflow}) & tentative\\
The strong shape effect needs self-propelled wall defects: with negative winding the $-\half$ wall layer survives elongation, and shielding weakens only modestly & ellipses at $k=-6$, $-12$, 2 seeds (Sec.~\ref{sec:smirror}) & tentative--good\\\bottomrule
\end{tabularx}
\end{table}
\clearpage
